In order to specify the hydrodynamic structure of an equilibrium configuration, one
generally gives---in addition to the ``gravitational potential,'' $\nu(r)$, and the mass inside
radius $r$, $m(r)$---also the total density of mass-energy, $\rho$, and the pressure, $p$, as functions
of $r$. The quantities $r$, $m$, $\rho$, and $p$ are related to each other by the mass equation,
\begin{equation}
  m = \int_0^r 4\pi r^2 \rho \, dr \;,
\label{eq:3a}
\end{equation}
the Tolman-Oppenheimer-Volkoff equation of hydrostatic equilibrium,
\begin{equation}
  \frac{dp}{dr} = -\frac{(\rho + p)(m + 4\pi r^3 p)}{r(r - 2m)} \;,
\label{eq:3b}
\end{equation}
the source equation for $\nu$,
\begin{equation}
  d\nu/dr = -2(\rho + p)^{-1}(dp/dr) \;,
\label{eq:3c}
\end{equation}
and an equation of state,
\begin{equation}
  p = p(\rho, Z_1, \ldots, Z_N) \;.
\label{eq:3d}
\end{equation}

Here $s$ is the entropy per baryon and $Z_1, \ldots, Z_N$ are the fractional abundances of the
various nuclear species. The adiabatic index, $\gamma \equiv \Gamma_1$ in the notation of RSSD---which
governs the response of the stellar material to adiabatic compressions is related to the
equation of state by
\begin{equation}
  \gamma = \bigl[(\rho + p)/p\bigr](\partial p/\partial \rho)_{s,\, Z_1,\, \ldots,\, Z_N} \;.
\label{eq:4}
\end{equation}

In order to construct a relativistic stellar model, one solves equations~(\ref{eq:3a}--\ref{eq:3d}),
together with certain equations of thermal structure and with gas characteristic equations
outlined in RSSD. (We do not write down here the remaining structure equations because
they play no role in the theory of pulsation.) Throughout the remainder of this
paper we assume that somebody has given us an equilibrium configuration in which
the structure parameters $r$, $m$, $\rho$, and $p$ satisfy equations~(\ref{eq:3a}--\ref{eq:3d}), and we examine the behavior of that configuration under non-radial perturbations.

\subsection*{(b) Fluid Displacement and Perturbation in the Geometry of Spacetime}

The small-amplitude motion of our perturbed configuration is described by the 3-vector
displacement, $\xi(t, r, \theta, \phi)$, of the fluid with respect to the equilibrium coordinate system $(t, r, \theta, \phi)$.
As a result of the fluid motion, the geometry of spacetime around and inside the equilibrium configuration is no longer described by the line element~(\ref{eq:1}). Rather, the
geometry fluctuates in a manner described by ten functions, $h_{\mu\nu} \equiv h_{\mu\nu}$, of $(t, r, \theta, \phi) = (x^0, x^1, x^2, x^3)$:
\begin{equation}
  ds^2 = (g_{\mu\nu}^{(0)} + h_{\mu\nu})\,dx^\mu\,dx^\nu \;.
\label{eq:5}
\end{equation}

The entire theory of non-radial pulsations consists of the study of the ``equations of
motion'' which govern the thirteen functions, $\xi_i$ and $h_{\mu\nu}$, of $(t, r, \theta, \phi)$.

In order to make the equations of motion as simple as possible, we analyze $\xi_i$ and
$h_{\mu\nu}$ into vectorial and tensorial spherical harmonics in a manner first employed by Regge
and Wheeler (1957). (See Appendix~A for details.) Each spherical harmonic is characterized
by the usual indices $l$ and $M$, and by a parity which can be $(-1)^l$ or
$(-1)^{l+1}$. Group-theoretic considerations---or, alternatively, an analysis of the equations
of motion---reveal that for small-amplitude motions there is no coupling between the
various harmonics $(l, m)$. Consequently, we can confine our attention to the various harmonics individually.

For small-amplitude motion in a given spherical harmonic, the equations of motion
